\section*{अध्याय \ref{ch:b3:radicals-carbenes} --- मूलक, कार्बीन और पुनर्विन्यास}

\begin{solution}{exo:b3:radicals-carbenes:1}
प्रोपेन में छह प्राथमिक और दो द्वितीयक हाइड्रोजन हैं। क्लोरीनन: $6 \times 1 : 2 \times 3.9 = 6 : 7.8$,
अर्थात् 43\,\% 1-क्लोरोप्रोपेन और 57\,\% 2-क्लोरोप्रोपेन। ब्रोमीनन: $6 : 164$, 4\,\% और
96\,\%।
\end{solution}

\begin{solution}{exo:b3:radicals-carbenes:2}
$\ce{CH2}$: त्रिक; $\ce{CCl2}$: एकल (क्लोरीन के एकाकी युग्म रिक्त p कक्षक को ऊपर उठाते हैं)।
डाइक्लोरोकार्बीन त्रिविमविशिष्ट रूप से जुड़ती है: \textit{trans}-1,1-डाइक्लोरो-2,3-डाइमेथिलसाइक्लोप्रोपेन।
\end{solution}

\begin{solution}{exo:b3:radicals-carbenes:3}
\textit{cis}-1,2-डाइएथिलसाइक्लोप्रोपेन: \omterm{def:b3:radicals-carbenes:carbenoid}{कार्बीनॉइड} एक ही चरण में एक फलक पर जुड़ता है।
\end{solution}

\begin{solution}{exo:b3:radicals-carbenes:4}
साइक्लोहेक्सेनोन: $\varepsilon$-कैप्रोलैक्टोन (ऑक्सेपेन-2-ओन)। ऐसीटोफ़ीनोन: फ़ेनिल ऐसीटेट,
$\ce{CH3CO2C6H5}$; मेथिल की तुलना में फ़ेनिल वरीयता से प्रवास करता है।
\end{solution}

\begin{solution}{exo:b3:radicals-carbenes:5}
नौ प्राथमिक और एक तृतीयक हाइड्रोजन। क्लोरीनन: $9 : 5.1$, 64\,\% 1-क्लोरो-2-मेथिलप्रोपेन
और 36\,\% 2-क्लोरो-2-मेथिलप्रोपेन। ब्रोमीनन: $9 : 1600$, 0.6\,\% और 99.4\,\%।
\end{solution}

\begin{solution}{exo:b3:radicals-carbenes:6}
$k_c = 3.0 \times k_H[\mathrm{Bu_3SnH}] = 3.0 \times \num{2.4e6} \times 0.020 = \qty{1.4e5}{s^{-1}}$।
\end{solution}

\begin{solution}{exo:b3:radicals-carbenes:7}
C5 पर आक्रमण: 5-exo-trig, अनुकूल; C6 पर: 6-endo-trig, नियमों द्वारा अनुकूल पर
यहाँ धीमा। 4-exo-tet: अनुकूल। 5-endo-trig: प्रतिकूल। 5-exo-dig: अनुकूल।
3-exo-dig: प्रतिकूल। 5-endo-dig: अनुकूल।
\end{solution}

\begin{solution}{exo:b3:radicals-carbenes:8}
जल का निकास एक तृतीयक, बेंज़िलिक धनायन देता है; पड़ोसी कार्बन पर एक फ़ेनिल
और एक मेथिल समूह खिसक सकते हैं, और उच्चतर क्षमता वाला फ़ेनिल प्रवास करता है।
ऑक्सीजन-स्थायीकृत धनायन एक प्रोटॉन खोता है: 3,3-डाइफ़ेनिलब्यूटेन-2-ओन।
\end{solution}

\begin{solution}{exo:b3:radicals-carbenes:9}\emergencystretch=3em
कर्टियस: $\ce{C6H5CON3 -> C6H5NCO + N2}$। जल के साथ कार्बैमिक अम्ल $\ce{CO2}$ खोता है: ऐनिलीन
(ऐनिलीन के और \omterm{def:b3:radicals-carbenes:curtius}{आइसोसायनेट} से मिलने पर कुछ डाइफ़ेनिलयूरिया बनता है)। एथेनॉल के साथ: एथिल
\textit{N}-फ़ेनिलकार्बामेट, $\ce{C6H5NHCO2C2H5}$।
\end{solution}

\begin{solution}{exo:b3:radicals-carbenes:10}
क्लोरीन: प्राथमिक $421.7 - 432 = \qty{-10}{kJ/mol}$, तृतीयक $405.4 - 432 = \qty{-27}{kJ/mol}$; ब्रोमीन: $+56$ और
$+39$। क्लोरीन के दोनों अपहरण ऊष्माक्षेपी हैं, प्रारंभिक संक्रमण अवस्थाओं के साथ जो
\qty{16}{kJ/mol} अंतर को कम अनुभव करती हैं; ब्रोमीन के दोनों अपहरण ऊष्माशोषी हैं,
विलंबित संक्रमण अवस्थाओं के साथ जिनकी सक्रियण ऊर्जाएँ इसके अधिकांश भाग से भिन्न होती हैं:
$\eu^{16000/(8.314 \times 298)} \approx 600$, प्रेक्षित अनुपात की कोटि।
\end{solution}

\begin{solution}{exo:b3:radicals-carbenes:11}
\textit{E} ऑक्सिम (OH, C2 के anti, जिस पर मेथिल समूह है): C2 प्रवास करता है, नाइट्रोजन
उसके बगल में प्रवेश करता है: 7-मेथिलऐज़ेपेन-2-ओन। \textit{Z} ऑक्सिम: C6 प्रवास करता है: 3-मेथिलऐज़ेपेन-2-ओन।
\omterm{def:b3:radicals-carbenes:beckmann}{बेकमान पुनर्विन्यास} में ऑक्सिम की ज्यामिति निर्णय करती है; \omterm{def:b3:radicals-carbenes:baeyer}{बायर--विलिगर
ऑक्सीकरण} में \omterm{def:b3:radicals-carbenes:migratory}{प्रवास क्षमता}, और द्वितीयक कार्बन खिसकता है:
7-मेथिलऑक्सेपेन-2-ओन।
\end{solution}

\begin{solution}{exo:b3:radicals-carbenes:12}
चक्रित अंश $k_c/(k_c + k_H c)$: 0.91, 0.49 और 0.09। 95\,\% के लिए: $k_Hc = k_c(1/0.95 - 1)$, अतः
$c = 0.0526 \times \num{2.3e5}/\num{2.4e6} = \qty{5.0e-3}{mol/L}$। इसे सिरिंज पंप से हाइड्राइड को धीरे मिलाकर, या हाइड्राइड को
पुनर्जनित करने वाले अपचायक के साथ उत्प्रेरकी मात्रा में टिन हैलाइड का उपयोग करके, बनाए रखा जाता है।
\end{solution}

\begin{solution}{pb:b3:radicals-carbenes:1}
\textbf{1.} $\ce{C6H10O + NH2OH -> C6H11NO + H2O}$।

\textbf{2.} योग मध्यवर्ती के हाइड्रॉक्सी समूह को जल के रूप में हटाने के लिए अम्ल
आवश्यक है, पर अधिक अम्ल हाइड्रॉक्सिलऐमीन को प्रोटॉनित कर देता है, जो फिर नहीं जुड़ती।

\textbf{3.} नहीं: C=N कार्बन से जुड़े दोनों कार्बन वलय की सममिति से
तुल्य हैं।

\textbf{4.} \textit{E} ऑक्सिम में OH मेथिल-धारी कार्बन C2 से दूर इंगित करता है (उसके
anti), \textit{Z} ऑक्सिम में उसकी ओर।

\textbf{5.} किसी विद्युतऋणी ऑक्सीजन वाले ऑक्सिम नाइट्रोजन पर प्रतिलोमन का अवरोध
ऊँचा होता है।

\textbf{6.} यह हाइड्रॉक्सी समूह को प्रोटॉनित (या सल्फ़ोनित) करके उसे अच्छा निर्गामी
समूह, जल, बना देता है।

\textbf{7.} निर्गामी समूह के anti स्थित वलय C--C आबंध कार्बन से नाइट्रोजन पर खिसकता है,
जो इस प्रकार वलय में निवेशित हो जाता है: छह परमाणु सात बन जाते हैं।

\textbf{8.} एक चक्रीय नाइट्रिलियम आयन।

\textbf{9.} $\ce{C6H11NO -> C6H11NO}$: एक समावयवन, ऑक्सिम से लैक्टम।

\emergencystretch=3em \textbf{10.} कैप्रोलैक्टम, ऐज़ेपेन-2-ओन, 6-ऐमीनोहेक्सेनोइक अम्ल का सप्तसदस्यीय चक्रीय
ऐमाइड; इसका वलय-विवृत बहुलकीकरण नायलॉन-6 देता है।

\textbf{11.} \qty{98.0}{g/mol} और \qty{113.0}{g/mol}।

\textbf{12.} $\qty{1.0e6}{g}/\qty{98.0}{g/mol} = \qty{10204}{mol}$; $\times 0.98 \times 0.98 = \qty{9800}{mol}$; $\times \qty{113.0}{g/mol} = \qty{1.11e6}{g}$, \qty{1.11}{t}।

\textbf{13.} ऑक्सिम: $\num{10204} \times 0.98 = \qty{1.00e4}{mol}$; $\ce{H2SO4}$: $\qty{1.30e4}{mol}$, जो उतने ही मोल
$\ce{(NH4)2SO4}$ देता है: $\qty{1.30e4}{mol} \times \qty{132.1}{g/mol} = \qty{1.7}{t}$, कैप्रोलैक्टम से अधिक।

\textbf{14.} $\ce{C6H10O + RCO3H -> C6H10O2 + RCOOH}$।

\textbf{15.} कार्बोनिल कार्बन पर पेरॉक्सी अम्ल का चतुष्फलकीय योगज (क्रीगी
मध्यवर्ती), एक पेरॉक्सी हेमीकीटल।

\textbf{16.} वलय का एक $\ce{CH2}$ समूह (दोनों तुल्य हैं) निकटतर पेरॉक्साइड
ऑक्सीजन पर खिसकता है; कार्बोक्सिलेट $\ce{RCOO-}$ निकलता है।

\emergencystretch=3em \textbf{17.} $\varepsilon$-कैप्रोलैक्टोन, ऑक्सेपेन-2-ओन; इसका वलय-विवृत बहुलकीकरण
पॉली($\varepsilon$-कैप्रोलैक्टोन) देता है, एक जैवनिम्नीकरणीय पॉलिएस्टर।

\textbf{18.} इसका एकमात्र उप-उत्पाद जल है, जबकि पेरॉक्सी अम्ल प्रति मोल लैक्टोन एक मोल
कार्बोक्सिलिक अम्ल छोड़ता है और स्वयं भंडारण के लिए ख़तरनाक है।

\textbf{19.} द्वितीयक कार्बन C2 प्रवास करता है: 7-मेथिलऑक्सेपेन-2-ओन।

\textbf{20.} हाँ: प्रवास करने वाला कार्बन अपना विन्यास बनाए रखता है।

\textbf{21.} 7-मेथिलऐज़ेपेन-2-ओन (C2 anti, प्रवास करता है)।

\textbf{22.} 3-मेथिलऐज़ेपेन-2-ओन (C6 प्रवास करता है)।

\textbf{23.} बेकमान: ऑक्सिम की ज्यामिति (anti समूह खिसकता है)।
बायर--विलिगर: \omterm{def:b3:radicals-carbenes:migratory}{प्रवास क्षमता} (अधिक प्रतिस्थापित कार्बन खिसकता है)।

\textbf{24.} इसके दोनों $\alpha$ कार्बन तुल्य हैं: एक ऑक्सिम, एक लैक्टम, एक लैक्टोन।

\textbf{25.} एक टन साइक्लोहेक्सेनोन लगभग \qty{1.11}{t} कैप्रोलैक्टम देता है।
\end{solution}
